% Revisão da ponderação por confiança — reexame de 17/08/2026.
% Fontes: Mestrado_PETR4/ponderacao_confianca.json
%         Mestrado_PETR4/ponderacao_confianca_apos17h.json
% Script: src/sentimento/testar_ponderacao_confianca.py
\subsection{Revisão da ponderação por confiança}
\label{sec:revisao_ponderacao}

A Seção~\ref{sec:refinamento_validacao} reporta um conjunto de resultados obtidos ao longo do
percurso de investigação, alguns dos quais precederam os instrumentos de avaliação hoje disponíveis.
Esta seção reexamina um deles --- a alegação de que a ponderação do índice de sentimento pela
confiança do classificador seria responsável pelo ganho preditivo --- e conclui pela sua revisão. O
reexame impôs-se por duas razões independentes.

A primeira é de coerência interna. A leitura original sustentava a preferência por
``um classificador que fornece um escore de confiança calibrado''. A
Seção~\ref{sec:achados_implementacao}, redigida posteriormente, estabeleceu que o escore publicado
pelo artefato \emph{não} é uma probabilidade calibrada: a declaração de
\texttt{problem\_type} como \texttt{multi\_label\_classification} faz com que a
biblioteca aplique a função sigmoide, e não a \textit{softmax} apropriada à classificação
exclusiva. As duas passagens eram, portanto, incompatíveis.

A segunda é de método. A medição original antecede os testes de associação com a volatilidade
introduzidos nas Seções~\ref{sec:relevancia_volatilidade} e~\ref{sec:auditoria_noticia}, e apoiava-se
exclusivamente na acurácia direcional --- eixo que, como se demonstrou, é insensível a melhorias do
componente textual.

\subsubsection{Confirmação da escala sigmoide}

Antes do reexame propriamente dito, cabe registrar uma verificação que dispensa qualquer suposição.
Em uma distribuição \textit{softmax} sobre três classes, a probabilidade da classe vencedora não
pode, por construção, situar-se abaixo de $1/3$: as três probabilidades somam a unidade, e a maior
delas é, no mínimo, a média. Observam-se, no corpus, \textbf{397 escores inferiores a $0{,}3333$},
com mínimo de $0{,}2845$ e máximo de $0{,}8729$. A escala não é, portanto,
\textit{softmax} --- constatação que corrobora, por via aritmética e independente, o achado da
Seção~\ref{sec:achados_implementacao}.

\subsubsection{O reexame}

Compararam-se quatro formas de agregar os mesmos rótulos em um índice diário: a ponderada pela
confiança, adotada na dissertação; a polaridade pura, que atribui $+1$, $0$ e $-1$ sem qualquer
peso; a restrição às notícias de confiança acima da mediana; e o saldo de votos. A avaliação
empregou tanto a associação com a volatilidade do pregão seguinte, apurada por \textit{bootstrap}
pareado, quanto a acurácia direcional sob o protocolo do Capítulo~\ref{chap:metodologia}. A
Tabela~\ref{tab:revisao_ponderacao} apresenta os resultados sobre o recorte de notícias posteriores
ao fechamento --- o mesmo empregado na medição original.

\begin{table}[htpb]
\centering
\caption{Reexame das construções do índice de sentimento (recorte posterior ao fechamento, 1.906 pregões).}
\label{tab:revisao_ponderacao}
\begin{tabular}{l c c c}
\hline
\textbf{Construção do índice} & \textbf{$|r|$ volatilidade} & \textbf{Acur.\ validação} & \textbf{Acur.\ teste} \\ \hline
Ponderada pela confiança & $0{,}0513$ & $56{,}64\%$ & $50{,}31\%$ \\
\textbf{Polaridade pura} & $0{,}0503$ & $53{,}85\%$ & $\mathbf{53{,}88\%}$ \\
Apenas alta confiança & $0{,}0436$ & $53{,}50\%$ & $53{,}04\%$ \\
Saldo de votos & $0{,}0503$ & $53{,}85\%$ & $53{,}88\%$ \\
\hline
\end{tabular}
\vspace{0.2cm}
{\small \\ Fonte: Elaborado pelo autor (\texttt{src/sentimento/testar\_ponderacao\_confianca.py}, 2026-08-17).}
\end{table}

Três observações decorrem da tabela, e as três contrariam a leitura original.

A primeira é que \textbf{a ponderação pela confiança piora a acurácia de teste}, e não a melhora:
$50{,}31\%$ contra $53{,}88\%$ da polaridade pura, uma diferença de $-3{,}57$ pontos percentuais. O
mesmo contraste sobre o corpus integral produz $-1{,}41$ ponto percentual, no mesmo sentido.

A segunda é que a variante ponderada exibe a maior acurácia de \emph{validação} ($56{,}64\%$) e a
menor de \emph{teste} ($50{,}31\%$) --- distância de mais de seis pontos percentuais que constitui a
assinatura característica da seleção sobre o conjunto de validação. É explicação plausível para a
cifra de $54{,}93\%$ originalmente reportada.

A terceira é que \textbf{a polaridade pura e o saldo de votos produzem resultados idênticos}, o que
não é coincidência: a média de $\{+1, 0, -1\}$ sobre as notícias do dia equivale, por identidade
algébrica, a $(n_{pos} - n_{neg})/n_{total}$. As duas construções são a mesma quantidade. A tabela
original, ao atribuir-lhes desempenhos distintos --- $54{,}53\%$ e $50{,}30\%$ ---, comparava
grandezas que não podiam divergir sob a definição declarada, discrepância cuja origem não foi
possível reconstituir, porquanto o código daquela suíte não foi preservado no repositório.

Sobre a associação com a volatilidade, a ponderação apresenta vantagem de $+0{,}0009$ no recorte
posterior ao fechamento, sem significância ($p = 0{,}789$), e de $+0{,}0084$ sobre o corpus integral,
esta significativa ($p = 0{,}021$). Trata-se, portanto, de um efeito real porém de magnitude
modesta: representa acréscimo de aproximadamente $6\%$ na intensidade da associação, contra os
$23\%$ obtidos pelo filtro de relevância da Seção~\ref{sec:relevancia_volatilidade}. A razão da
modéstia é aritmética: as séries ponderada e não ponderada correlacionam-se a $0{,}9903$, de sorte
que a ponderação redistribui muito pouco.

\subsubsection{Consequências}

A conclusão revista pode ser enunciada em três proposições.

A tese de que ``o sinal preditivo reside na intensidade e no grau de certeza do modelo''
\textbf{não se sustenta}. A ponderação pela confiança contribui, quando muito, com um acréscimo
marginal à associação com a volatilidade, e prejudica a acurácia direcional. Mantém-se a construção
ponderada na dissertação por continuidade com os resultados já reportados, mas sem lhe atribuir o
mérito que a leitura original lhe conferia.

Segue-se, como corolário favorável, que \textbf{a escala equivocada do escore de confiança tem
consequência prática desprezível} sobre os resultados agregados. Recalcular o índice com a função
\textit{softmax} apropriada --- procedimento que exigiria nova inferência sobre as 205.697 notícias,
com processamento em unidade de processamento gráfico --- alteraria uma ponderação cujo efeito é, ele
próprio, marginal. O achado da Seção~\ref{sec:achados_implementacao} permanece relevante como
alerta de implementação para terceiros, mas não compromete as conclusões aqui reportadas.

Reforça-se, por fim, o padrão já identificado na Seção~\ref{sec:relevancia_volatilidade}: entre as
intervenções testadas, aquelas que operam sobre a \emph{seleção} das notícias produzem ganhos de
magnitude superior às que operam sobre a \emph{forma de agregá-las} ou sobre a leitura que o modelo
faz de cada uma. O filtro de relevância acrescenta $23\%$ à associação; a ponderação por confiança,
$6\%$; as oito intervenções sobre o classificador, nada.
